Soit $R$ et $S$ des points distincts appartenant à un cercle $\Omega$ tels que le segment $[RS]$ n'est pas un
diamètre de $\Omega$. Soit $\ell$ la tangente à $\Omega$ en $R$. Le point $T$ est tel que $S$ est le milieu du
segment $[RT]$. Le point $J$ est choisi sur le plus petit arc $\overset{\frown}{RS}$ de $\Omega$ de sorte que
le cercle $\Gamma$ circonscrit au triangle $JST$ rencontre $\ell$ en deux points distincts. Soit $A$ le point
commun de $\Gamma$ et $\ell$ qui est le plus proche de $R$. La droite $(AJ)$ recoupe $\Omega$ en $K$. Prouver
que la droite $(KT)$ est tangente à $\Gamma$.
